\section{Normal cofinal maps dominating a function}

\import{function.tex}
\import{ordinal.tex}
\import{cardinal/infinite.tex}
\import{cardinal/normal-map.tex}
\import{cardinal/cofinality.tex}
\import{recursion/transfinite.tex}

\subsection{The total recursion step}

\begin{definition}\label{normal_dominating_next}
    $\normalDominatingNext{t}{\beta}{g}
        =\suc{(\beta\union g(\unions{\beta}))
            \union t(g(\unions{\beta}))}$.
\end{definition}

\begin{definition}\label{normal_dominating_successor_guard}
    $\normalDominatingSuccessorGuard{t}{\lambda}{\beta}{g}$ iff
        $\beta$ is a successor ordinal and
        $\normalDominatingNext{t}{\beta}{g}\in\lambda$.
\end{definition}

\begin{definition}\label{normal_dominating_limit_guard}
    $\normalDominatingLimitGuard{\lambda}{\beta}{g}$ iff
        $\beta$ is a limit ordinal and
        $\unions{\ran{g}}\in\lambda$.
\end{definition}

\begin{definition}\label{normal_dominating_step}
    $\normalDominatingStep{t}{\lambda}
        =\{w\in
            ((\lambda\times\pow{\lambda\times\lambda})\times\lambda)
            \mid\text{
            ($\fst{\fst{w}}=\emptyset$ and $\snd{w}=\emptyset$) or
            ($\normalDominatingSuccessorGuard{t}{\lambda}
                {\fst{\fst{w}}}{\snd{\fst{w}}}$ and
                $\snd{w}=\normalDominatingNext{t}
                    {\fst{\fst{w}}}{\snd{\fst{w}}}$) or
            ($\normalDominatingLimitGuard{\lambda}
                {\fst{\fst{w}}}{\snd{\fst{w}}}$ and
                $\snd{w}=\unions{\ran{\snd{\fst{w}}}}$) or
            ($\fst{\fst{w}}\neq\emptyset$ and
                not $\normalDominatingSuccessorGuard{t}{\lambda}
                    {\fst{\fst{w}}}{\snd{\fst{w}}}$ and
                not $\normalDominatingLimitGuard{\lambda}
                    {\fst{\fst{w}}}{\snd{\fst{w}}}$ and
                $\snd{w}=\emptyset$)}\}$.
\end{definition}

\begin{proposition}\label{normal_dominating_step_iff}
    Let $q\in\lambda\times\pow{\lambda\times\lambda}$ and
        $y\in\lambda$.
    Then $(q,y)\in\normalDominatingStep{t}{\lambda}$ iff
        ($\fst{q}=\emptyset$ and $y=\emptyset$) or
        ($\normalDominatingSuccessorGuard{t}{\lambda}
            {\fst{q}}{\snd{q}}$ and
            $y=\normalDominatingNext{t}{\fst{q}}{\snd{q}}$) or
        ($\normalDominatingLimitGuard{\lambda}{\fst{q}}{\snd{q}}$ and
            $y=\unions{\ran{\snd{q}}}$) or
        ($\fst{q}\neq\emptyset$ and
            not $\normalDominatingSuccessorGuard{t}{\lambda}
                {\fst{q}}{\snd{q}}$ and
            not $\normalDominatingLimitGuard{\lambda}
                {\fst{q}}{\snd{q}}$ and
            $y=\emptyset$).
\end{proposition}
\begin{proof}
    $(q,y)\in
        (\lambda\times\pow{\lambda\times\lambda})\times\lambda$ by
        \obv{times_tuple_intro}.
    Follows by \auto{normal_dominating_step,fst_eq,snd_eq}.
\end{proof}

\begin{theorem}\label{normal_dominating_step_type}
    Suppose $\emptyset\in\lambda$.
    Then $\normalDominatingStep{t}{\lambda}
        \in\funs{\lambda\times\pow{\lambda\times\lambda}}{\lambda}$.
\end{theorem}
\begin{proof}
    Let $D=\lambda\times\pow{\lambda\times\lambda}$.
    Let $T=\normalDominatingStep{t}{\lambda}$.
    $T\subseteq D\times\lambda$ by
        \auto{normal_dominating_step,subseteq}.
    Hence $T$ is a relation by
        \auto{relation,subseteq,times_elem_is_tuple}.
    Show that $T$ is right-unique.
    \begin{subproof}
        Show that for all $q,y,z$ such that $(q,y),(q,z)\in T$
            we have $y=z$.
        \begin{subproof}
            Fix $q,y,z$.
            Assume $(q,y),(q,z)\in T$.
            Hence $q\in D$ and $y,z\in\lambda$ by
                \obv{subseteq,times_tuple_elim}.
            Hence
                ($\fst{q}=\emptyset$ and $y=\emptyset$) or
                ($\normalDominatingSuccessorGuard{t}{\lambda}
                    {\fst{q}}{\snd{q}}$ and
                    $y=\normalDominatingNext{t}{\fst{q}}{\snd{q}}$) or
                ($\normalDominatingLimitGuard{\lambda}
                    {\fst{q}}{\snd{q}}$ and
                    $y=\unions{\ran{\snd{q}}}$) or
                ($\fst{q}\neq\emptyset$ and
                    not $\normalDominatingSuccessorGuard{t}{\lambda}
                        {\fst{q}}{\snd{q}}$ and
                    not $\normalDominatingLimitGuard{\lambda}
                        {\fst{q}}{\snd{q}}$ and
                    $y=\emptyset$) by
                \auto{normal_dominating_step_iff}.
            Hence
                ($\fst{q}=\emptyset$ and $z=\emptyset$) or
                ($\normalDominatingSuccessorGuard{t}{\lambda}
                    {\fst{q}}{\snd{q}}$ and
                    $z=\normalDominatingNext{t}{\fst{q}}{\snd{q}}$) or
                ($\normalDominatingLimitGuard{\lambda}
                    {\fst{q}}{\snd{q}}$ and
                    $z=\unions{\ran{\snd{q}}}$) or
                ($\fst{q}\neq\emptyset$ and
                    not $\normalDominatingSuccessorGuard{t}{\lambda}
                        {\fst{q}}{\snd{q}}$ and
                    not $\normalDominatingLimitGuard{\lambda}
                        {\fst{q}}{\snd{q}}$ and
                    $z=\emptyset$) by
                \auto{normal_dominating_step_iff}.
            Follows by
                \auto{normal_dominating_successor_guard,
                    normal_dominating_limit_guard,
                    zero_not_successorordinal,zero_not_limitordinal,
                    limit_ordinal_is_not_successor_ordinal}.
        \end{subproof}
        Follows by \obv{rightunique}.
    \end{subproof}
    Show that $\dom{T}=D$.
    \begin{subproof}
        Show that $\dom{T}\subseteq D$.
        \begin{subproof}
            Follows by \auto{subseteq,dom_iff,times_tuple_elim}.
        \end{subproof}
        Show that $D\subseteq\dom{T}$.
        \begin{subproof}
            Show that every element of $D$ is an element of $\dom{T}$.
            \begin{subproof}
                Fix $q\in D$.
                \begin{byCase}
                    \caseOf{$\fst{q}=\emptyset$.}
                        $(q,\emptyset)\in T$ by
                            \auto{normal_dominating_step_iff}.
                        Follows by \auto{dom_iff}.
                    \caseOf{$\fst{q}\neq\emptyset$.}
                        \begin{byCase}
                            \caseOf{$\normalDominatingSuccessorGuard{t}
                                    {\lambda}{\fst{q}}{\snd{q}}$.}
                                $\normalDominatingNext{t}{\fst{q}}{\snd{q}}
                                    \in\lambda$ by
                                    \obv{normal_dominating_successor_guard}.
                                Hence $(q,\normalDominatingNext{t}
                                        {\fst{q}}{\snd{q}})\in T$ by
                                    \auto{normal_dominating_step_iff}.
                                Follows by \auto{dom_iff}.
                            \caseOf{not
                                    $\normalDominatingSuccessorGuard{t}
                                        {\lambda}{\fst{q}}{\snd{q}}$.}
                                \begin{byCase}
                                    \caseOf{$\normalDominatingLimitGuard
                                            {\lambda}{\fst{q}}{\snd{q}}$.}
                                        $\unions{\ran{\snd{q}}}\in\lambda$ by
                                            \obv{normal_dominating_limit_guard}.
                                        Hence $(q,\unions{\ran{\snd{q}}})
                                            \in T$ by
                                            \auto{normal_dominating_step_iff}.
                                        Follows by \auto{dom_iff}.
                                    \caseOf{not
                                            $\normalDominatingLimitGuard
                                                {\lambda}{\fst{q}}{\snd{q}}$.}
                                        $(q,\emptyset)\in T$ by
                                            \auto{normal_dominating_step_iff}.
                                        Follows by \auto{dom_iff}.
                                \end{byCase}
                        \end{byCase}
                \end{byCase}
            \end{subproof}
            Follows by \obv{subseteq}.
        \end{subproof}
        Follows by \obv{subseteq_antisymmetric}.
    \end{subproof}
    Show that for every $q\in\dom{T}$ we have $T(q)\in\lambda$.
    \begin{subproof}
        Fix $q\in\dom{T}$.
        $(q,T(q))\in T$ by \obv{function_apply_elim}.
        Hence $(q,T(q))\in D\times\lambda$ by \obv{subseteq}.
        Follows by \obv{times_tuple_elim}.
    \end{subproof}
    Hence $T$ is a function to $\lambda$.
    Follows by \auto{funs_intro}.
\end{proof}

\begin{proposition}\label{normal_dominating_step_zero}
    Suppose $\emptyset\in\lambda$.
    Let $g\in\pow{\lambda\times\lambda}$.
    Then $(\normalDominatingStep{t}{\lambda})((\emptyset,g))
        =\emptyset$.
\end{proposition}
\begin{proof}
    Let $q=(\emptyset,g)$.
    $q\in\lambda\times\pow{\lambda\times\lambda}$ by
        \obv{times_tuple_intro}.
    $\normalDominatingStep{t}{\lambda}
        \in\funs{\lambda\times\pow{\lambda\times\lambda}}{\lambda}$ by
        \obv{normal_dominating_step_type}.
    Hence $\normalDominatingStep{t}{\lambda}$ is a function by
        \obv{funs_is_function}.
    $(q,\emptyset)\in\normalDominatingStep{t}{\lambda}$ by
        \auto{normal_dominating_step_iff,fst_eq,snd_eq}.
    Hence $(\normalDominatingStep{t}{\lambda})(q)=\emptyset$ by
        \obv{function_apply_intro}.
\end{proof}

\begin{proposition}\label{normal_dominating_step_successor}
    Suppose $\emptyset\in\lambda$.
    Let $\beta\in\lambda$ and
        $g\in\pow{\lambda\times\lambda}$.
    Suppose $\normalDominatingSuccessorGuard{t}{\lambda}{\beta}{g}$.
    Then $(\normalDominatingStep{t}{\lambda})((\beta,g))
        =\normalDominatingNext{t}{\beta}{g}$.
\end{proposition}
\begin{proof}
    Let $q=(\beta,g)$.
    $q\in\lambda\times\pow{\lambda\times\lambda}$ by
        \obv{times_tuple_intro}.
    $\normalDominatingNext{t}{\beta}{g}\in\lambda$ by
        \obv{normal_dominating_successor_guard}.
    $\normalDominatingStep{t}{\lambda}
        \in\funs{\lambda\times\pow{\lambda\times\lambda}}{\lambda}$ by
        \obv{normal_dominating_step_type}.
    Hence $\normalDominatingStep{t}{\lambda}$ is a function by
        \obv{funs_is_function}.
    $(q,\normalDominatingNext{t}{\beta}{g})
        \in\normalDominatingStep{t}{\lambda}$ by
        \auto{normal_dominating_step_iff,fst_eq,snd_eq}.
    Hence $(\normalDominatingStep{t}{\lambda})(q)
        =\normalDominatingNext{t}{\beta}{g}$ by
        \obv{function_apply_intro}.
\end{proof}

\begin{proposition}\label{normal_dominating_step_limit}
    Suppose $\emptyset\in\lambda$.
    Let $\beta\in\lambda$ and
        $g\in\pow{\lambda\times\lambda}$.
    Suppose $\normalDominatingLimitGuard{\lambda}{\beta}{g}$.
    Then $(\normalDominatingStep{t}{\lambda})((\beta,g))
        =\unions{\ran{g}}$.
\end{proposition}
\begin{proof}
    Let $q=(\beta,g)$.
    $q\in\lambda\times\pow{\lambda\times\lambda}$ by
        \obv{times_tuple_intro}.
    $\unions{\ran{g}}\in\lambda$ by
        \obv{normal_dominating_limit_guard}.
    $\normalDominatingStep{t}{\lambda}
        \in\funs{\lambda\times\pow{\lambda\times\lambda}}{\lambda}$ by
        \obv{normal_dominating_step_type}.
    Hence $\normalDominatingStep{t}{\lambda}$ is a function by
        \obv{funs_is_function}.
    $(q,\unions{\ran{g}})\in\normalDominatingStep{t}{\lambda}$ by
        \auto{normal_dominating_step_iff,fst_eq,snd_eq}.
    Hence $(\normalDominatingStep{t}{\lambda})(q)
        =\unions{\ran{g}}$ by \obv{function_apply_intro}.
\end{proof}

\begin{proposition}\label{normal_dominating_step_fallback}
    Suppose $\emptyset\in\lambda$.
    Let $\beta\in\lambda$ and
        $g\in\pow{\lambda\times\lambda}$.
    Suppose $\beta\neq\emptyset$ and
        not $\normalDominatingSuccessorGuard{t}{\lambda}{\beta}{g}$ and
        not $\normalDominatingLimitGuard{\lambda}{\beta}{g}$.
    Then $(\normalDominatingStep{t}{\lambda})((\beta,g))
        =\emptyset$.
\end{proposition}
\begin{proof}
    Let $q=(\beta,g)$.
    $q\in\lambda\times\pow{\lambda\times\lambda}$ by
        \obv{times_tuple_intro}.
    $\normalDominatingStep{t}{\lambda}
        \in\funs{\lambda\times\pow{\lambda\times\lambda}}{\lambda}$ by
        \obv{normal_dominating_step_type}.
    Hence $\normalDominatingStep{t}{\lambda}$ is a function by
        \obv{funs_is_function}.
    $(q,\emptyset)\in\normalDominatingStep{t}{\lambda}$ by
        \auto{normal_dominating_step_iff,fst_eq,snd_eq}.
    Hence $(\normalDominatingStep{t}{\lambda})(q)=\emptyset$ by
        \obv{function_apply_intro}.
\end{proof}

\subsection{The actual recursion}

\begin{proposition}\label{normal_dominating_recursion_type}
    Let $\lambda$ be a regular cardinal.
    Let $t\in\funs{\lambda}{\lambda}$.
    Then
    \[
        \ordinalRecursion{\lambda}{\lambda}
            {\normalDominatingStep{t}{\lambda}}
        \in\funs{\lambda}{\lambda}
    \].
\end{proposition}
\begin{proof}
    Hence $\lambda$ is a limit ordinal by
        \obv{infinite_cardinal_is_limit_ordinal}.
    Hence $\lambda$ is an ordinal and $\emptyset\in\lambda$ by
        \obv{limit_ordinal}.
    $\normalDominatingStep{t}{\lambda}
        \in\funs{\lambda\times\pow{\lambda\times\lambda}}{\lambda}$ by
        \obv{normal_dominating_step_type}.
    Follows by \obv{ordinal_recursion_type}.
\end{proof}

\begin{proposition}\label{normal_dominating_recursion_equation}
    Let $\lambda$ be a regular cardinal.
    Let $t\in\funs{\lambda}{\lambda}$.
    Let $\beta\in\lambda$.
    Then
    \[
        (\ordinalRecursion{\lambda}{\lambda}
            {\normalDominatingStep{t}{\lambda}})(\beta)
        =(\normalDominatingStep{t}{\lambda})
            ((\beta,\restrl{\ordinalRecursion{\lambda}{\lambda}
                {\normalDominatingStep{t}{\lambda}}}{\beta}))
    \].
\end{proposition}
\begin{proof}
    Hence $\lambda$ is a limit ordinal by
        \obv{infinite_cardinal_is_limit_ordinal}.
    Hence $\lambda$ is an ordinal and $\emptyset\in\lambda$ by
        \obv{limit_ordinal}.
    $\normalDominatingStep{t}{\lambda}
        \in\funs{\lambda\times\pow{\lambda\times\lambda}}{\lambda}$ by
        \obv{normal_dominating_step_type}.
    Follows by \obv{ordinal_recursion_equation}.
\end{proof}

\begin{proposition}\label{normal_dominating_recursion_zero}
    Let $\lambda$ be a regular cardinal.
    Let $t\in\funs{\lambda}{\lambda}$.
    Then
    \[
        (\ordinalRecursion{\lambda}{\lambda}
            {\normalDominatingStep{t}{\lambda}})(\emptyset)
        =\emptyset
    \].
\end{proposition}
\begin{proof}
    Hence $\lambda$ is a limit ordinal by
        \obv{infinite_cardinal_is_limit_ordinal}.
    Hence $\emptyset\in\lambda$ by \obv{limit_ordinal}.
    $\ordinalRecursion{\lambda}{\lambda}
            {\normalDominatingStep{t}{\lambda}}
        \in\funs{\lambda}{\lambda}$ by
        \obv{normal_dominating_recursion_type}.
    $\restrl{\ordinalRecursion{\lambda}{\lambda}
        {\normalDominatingStep{t}{\lambda}}}{\emptyset}=\emptyset$ by
        \auto{setext,restrl,emptyset}.
    $(\ordinalRecursion{\lambda}{\lambda}
            {\normalDominatingStep{t}{\lambda}})(\emptyset)
        =(\normalDominatingStep{t}{\lambda})
            ((\emptyset,\restrl{\ordinalRecursion{\lambda}{\lambda}
                {\normalDominatingStep{t}{\lambda}}}{\emptyset}))$ by
        \obv{normal_dominating_recursion_equation}.
    $(\normalDominatingStep{t}{\lambda})((\emptyset,\emptyset))
        =\emptyset$ by
        \obv{normal_dominating_step_zero,powerset_bottom}.
\end{proof}

\begin{theorem}\label{normal_dominating_recursion_successor}
    Let $\lambda$ be a regular cardinal.
    Let $t\in\funs{\lambda}{\lambda}$.
    Let $\mu\in\lambda$.
    Then
    \[
        (\ordinalRecursion{\lambda}{\lambda}
            {\normalDominatingStep{t}{\lambda}})(\suc{\mu})
        =\normalDominatingNext{t}{\suc{\mu}}
            {\restrl{\ordinalRecursion{\lambda}{\lambda}
                {\normalDominatingStep{t}{\lambda}}}{\suc{\mu}}}
    \]
    and
    \[
        t((\ordinalRecursion{\lambda}{\lambda}
            {\normalDominatingStep{t}{\lambda}})(\mu))
        \in
        (\ordinalRecursion{\lambda}{\lambda}
            {\normalDominatingStep{t}{\lambda}})(\suc{\mu})
    \].
\end{theorem}
\begin{proof}
    $\lambda$ is an infinite cardinal and
        $\cofinality{\lambda}=\lambda$ by \obv{regular_cardinal}.
    Hence $\lambda$ is a limit ordinal by
        \obv{infinite_cardinal_is_limit_ordinal}.
    Hence $\lambda$ is an ordinal and $\emptyset\in\lambda$ by
        \obv{limit_ordinal}.
    $\suc{\mu}\in\lambda$ by \obv{suc_elem_limitordinal}.
    $\ordinalRecursion{\lambda}{\lambda}
            {\normalDominatingStep{t}{\lambda}}
        \in\funs{\lambda}{\lambda}$ by
        \obv{normal_dominating_recursion_type}.
    Let $q=\ordinalRecursion{\lambda}{\lambda}
        {\normalDominatingStep{t}{\lambda}}$.
    $q\in\funs{\lambda}{\lambda}$.
    Hence $q$ is a function and $\dom{q}=\lambda$ by
        \obv{funs_is_function,funs}.
    $q(\mu)\in\lambda$ by \obv{funs_type_apply}.
    Hence $t(q(\mu))\in\lambda$ by \obv{funs_type_apply}.
    $\suc{\mu}$, $q(\mu)$, and $t(q(\mu))$ are ordinals by
        \obv{elem_of_ordinal_is_ordinal}.
    $(\suc{\mu}\union q(\mu))\union t(q(\mu))\in\lambda$ by
        \obv{three_ordinal_union_below}.
    Let $A=(\suc{\mu}\union q(\mu))\union t(q(\mu))$.
    $A\in\lambda$ and $A$ is an ordinal by
        \obv{elem_of_ordinal_is_ordinal}.
    Hence $\suc{A}\in\lambda$ by \obv{suc_elem_limitordinal}.
    $\suc{\mu}$ is a successor ordinal by
        \obv{ordinal_successor_data,elem_of_ordinal_is_ordinal}.
    $\suc{\mu}\subseteq\lambda$ by
        \obv{elem_of_ordinal_is_ordinal,ordinal_elem_implies_subseteq}.
    Hence $\suc{\mu}\subseteq\dom{q}$.
    $\mu\in\suc{\mu}$ by \obv{suc_intro_self}.
    $\unions{\suc{\mu}}=\mu$ by
        \obv{ordinal_successor_data,elem_of_ordinal_is_ordinal}.
    $(\restrl{q}{\suc{\mu}})(\mu)=q(\mu)$ by
        \obv{restrl_of_function_apply}.
    $\normalDominatingNext{t}{\suc{\mu}}{\restrl{q}{\suc{\mu}}}
        =\suc{A}$ by
        \obv{normal_dominating_next}.
    Hence $\normalDominatingSuccessorGuard{t}{\lambda}{\suc{\mu}}
        {\restrl{q}{\suc{\mu}}}$ by
        \obv{normal_dominating_successor_guard}.
    $q\in\rels{\lambda}{\lambda}$ by \obv{funs}.
    Hence $q\subseteq\lambda\times\lambda$ by \obv{rels_elim}.
    $\restrl{q}{\suc{\mu}}\subseteq q$ by \obv{restrl_subseteq}.
    Hence $\restrl{q}{\suc{\mu}}\subseteq\lambda\times\lambda$ by
        \obv{subseteq_transitive}.
    Hence $\restrl{q}{\suc{\mu}}\in\pow{\lambda\times\lambda}$ by
        \obv{powerset_intro}.
    $q(\suc{\mu})=(\normalDominatingStep{t}{\lambda})
        ((\suc{\mu},\restrl{q}{\suc{\mu}}))$ by
        \auto{normal_dominating_recursion_equation}.
    Hence $q(\suc{\mu})=
        \normalDominatingNext{t}{\suc{\mu}}{\restrl{q}{\suc{\mu}}}$ by
        \obv{normal_dominating_step_successor}.
    $t(q(\mu))\subseteq A$ by
        \obv{union_upper_right}.
    $t(q(\mu))\in A$ or $t(q(\mu))=A$ by
        \obv{ordinal_subseteq_implies_elem_or_eq}.
    Hence $t(q(\mu))\in\suc{A}$ by \auto{suc_iff}.
    Hence $t(q(\mu))\in q(\suc{\mu})$ by
        \obv{normal_dominating_next}.
\end{proof}

\begin{proposition}\label{normal_dominating_recursion_limit}
    Let $\lambda$ be a regular cardinal.
    Let $t\in\funs{\lambda}{\lambda}$.
    Let $\beta\in\lambda$.
    Suppose $\beta$ is a limit ordinal.
    Then
    \[
        (\ordinalRecursion{\lambda}{\lambda}
            {\normalDominatingStep{t}{\lambda}})(\beta)
        =\unions{\ran{\restrl{\ordinalRecursion{\lambda}{\lambda}
            {\normalDominatingStep{t}{\lambda}}}{\beta}}}
    \].
\end{proposition}
\begin{proof}
    $\lambda$ is an infinite cardinal and
        $\cofinality{\lambda}=\lambda$ by \obv{regular_cardinal}.
    Hence $\lambda$ is a limit ordinal by
        \obv{infinite_cardinal_is_limit_ordinal}.
    Hence $\lambda$ is an ordinal and $\emptyset\in\lambda$ by
        \obv{limit_ordinal}.
    $\ordinalRecursion{\lambda}{\lambda}
            {\normalDominatingStep{t}{\lambda}}
        \in\funs{\lambda}{\lambda}$ by
        \obv{normal_dominating_recursion_type}.
    Let $q=\ordinalRecursion{\lambda}{\lambda}
        {\normalDominatingStep{t}{\lambda}}$.
    $q\in\funs{\lambda}{\lambda}$.
    Hence $q$ is a function and $\dom{q}=\lambda$ by
        \obv{funs_is_function,funs}.
    $\beta$ is an ordinal and $\beta\subseteq\lambda$ by
        \obv{limit_ordinal,ordinal_elem_implies_subseteq}.
    Hence $\restrl{q}{\beta}\in\funs{\beta}{\lambda}$ by
        \obv{funs_restrl}.
    $\beta\in\cofinality{\lambda}$.
    Hence $\unions{\ran{\restrl{q}{\beta}}}\in\lambda$ by
        \obv{small_domain_range_union_below}.
    Hence $\normalDominatingLimitGuard{\lambda}{\beta}{\restrl{q}{\beta}}$
        by \obv{normal_dominating_limit_guard}.
    $q\in\rels{\lambda}{\lambda}$ by \obv{funs}.
    Hence $q\subseteq\lambda\times\lambda$ by \obv{rels_elim}.
    $\restrl{q}{\beta}\subseteq q$ by \obv{restrl_subseteq}.
    Hence $\restrl{q}{\beta}\subseteq\lambda\times\lambda$ by
        \obv{subseteq_transitive}.
    Hence $\restrl{q}{\beta}\in\pow{\lambda\times\lambda}$ by
        \obv{powerset_intro}.
    $q(\beta)=(\normalDominatingStep{t}{\lambda})
        ((\beta,\restrl{q}{\beta}))$ by
        \auto{normal_dominating_recursion_equation}.
    Hence $q(\beta)=\unions{\ran{\restrl{q}{\beta}}}$ by
        \auto{normal_dominating_step_limit}.
\end{proof}

\subsection{Normality, cofinality, and domination}

\begin{theorem}\label{normal_dominating_recursion_growth}
    Let $\lambda$ be a regular cardinal.
    Let $t\in\funs{\lambda}{\lambda}$.
    Let $\beta\in\lambda$.
    Then
    \[
        \beta\subseteq
        (\ordinalRecursion{\lambda}{\lambda}
            {\normalDominatingStep{t}{\lambda}})(\beta)
    \]
    and for every $\gamma\in\beta$ we have
    \[
        (\ordinalRecursion{\lambda}{\lambda}
            {\normalDominatingStep{t}{\lambda}})(\gamma)
        \in
        (\ordinalRecursion{\lambda}{\lambda}
            {\normalDominatingStep{t}{\lambda}})(\beta)
    \].
\end{theorem}
\begin{proof}[Proof by \in-induction on $\beta$]
    Hence $\lambda$ is a limit ordinal by
        \obv{infinite_cardinal_is_limit_ordinal}.
    Hence $\lambda$ is an ordinal by \obv{limit_ordinal}.
    $\ordinalRecursion{\lambda}{\lambda}
            {\normalDominatingStep{t}{\lambda}}
        \in\funs{\lambda}{\lambda}$ by
        \obv{normal_dominating_recursion_type}.
    $\beta$ is an ordinal by \obv{elem_of_ordinal_is_ordinal}.
    For every $\delta\in\beta$ we have $\delta\in\lambda$ by
        \obv{ordinal_transitivity}.
    Hence for every $\delta\in\beta$ we have
        $\delta\subseteq
            (\ordinalRecursion{\lambda}{\lambda}
                {\normalDominatingStep{t}{\lambda}})(\delta)$
        and for every $\xi\in\delta$ we have
            $(\ordinalRecursion{\lambda}{\lambda}
                {\normalDominatingStep{t}{\lambda}})(\xi)
            \in
            (\ordinalRecursion{\lambda}{\lambda}
                {\normalDominatingStep{t}{\lambda}})(\delta)$.
    Let $q=\ordinalRecursion{\lambda}{\lambda}
        {\normalDominatingStep{t}{\lambda}}$.
    $q\in\funs{\lambda}{\lambda}$.
    Hence $q$ is a function and $\dom{q}=\lambda$ by
        \obv{funs_is_function,funs}.
    $q(\beta)\in\lambda$ by \obv{funs_type_apply}.
    Hence for every $\delta\in\beta$ we have
        $\delta\subseteq q(\delta)$ and for every $\xi\in\delta$ we have
            $q(\xi)\in q(\delta)$ by \auto.
    \begin{byCase}
        \caseOf{$\beta=\emptyset$.}
            $\beta\subseteq q(\beta)$ by \obv{emptyset_subseteq}.
            For every $\gamma\in\beta$ we have
                $q(\gamma)\in q(\beta)$ by \auto{emptyset}.
            Follows by \auto.
        \caseOf{$\beta\neq\emptyset$.}
            $\emptyset\in\beta$ by
                \obv{ordinal_empty_or_emptyset_elem}.
            Hence $\emptyset\precedes\beta$.
            Hence $\beta$ is a limit ordinal or
                $\beta$ is a successor ordinal by
                \obv{positive_ordinal_is_limit_or_successor}.
            \begin{byCase}
                \caseOf{$\beta$ is a limit ordinal.}
                    $q(\beta)=\unions{\ran{\restrl{q}{\beta}}}$ by
                        \auto{normal_dominating_recursion_limit}.
                    $\beta\subseteq\lambda$ by
                        \obv{ordinal_elem_implies_subseteq}.
                    Hence $\restrl{q}{\beta}\in\funs{\beta}{\lambda}$ by
                        \obv{funs_restrl}.
                    Hence $\restrl{q}{\beta}$ is a function and
                        $\dom{\restrl{q}{\beta}}=\beta$ by
                        \obv{funs_is_function,funs}.
                    Show that for every $\gamma\in\beta$ we have
                        $\gamma\in q(\beta)$.
                    \begin{subproof}
                        Fix $\gamma\in\beta$.
                        $\suc{\gamma}\in\beta$ by
                            \obv{suc_elem_limitordinal}.
                        Hence $\suc{\gamma}\subseteq q(\suc{\gamma})$.
                        $\gamma\in\suc{\gamma}$ by \obv{suc_intro_self}.
                        Hence $\gamma\in q(\suc{\gamma})$ by
                            \obv{subseteq}.
                        $(\restrl{q}{\beta})(\suc{\gamma})
                            =q(\suc{\gamma})$ by
                            \obv{restrl_of_function_apply,
                                ordinal_elem_implies_subseteq}.
                        Hence $q(\suc{\gamma})
                            \in\ran{\restrl{q}{\beta}}$ by
                            \auto{fun_ran_iff}.
                        Hence $\gamma\in
                            \unions{\ran{\restrl{q}{\beta}}}$ by
                            \auto{unions_iff}.
                    \end{subproof}
                    Hence $\beta\subseteq q(\beta)$ by \obv{subseteq}.
                    Show that for every $\gamma\in\beta$ we have
                        $q(\gamma)\in q(\beta)$.
                    \begin{subproof}
                        Fix $\gamma\in\beta$.
                        $\suc{\gamma}\in\beta$ by
                            \obv{suc_elem_limitordinal}.
                        For every $\xi\in\suc{\gamma}$ we have
                            $q(\xi)\in q(\suc{\gamma})$ by \auto.
                        $\gamma\in\suc{\gamma}$ by \obv{suc_intro_self}.
                        Hence $q(\gamma)\in q(\suc{\gamma})$.
                        $(\restrl{q}{\beta})(\suc{\gamma})
                            =q(\suc{\gamma})$ by
                            \obv{restrl_of_function_apply,
                                ordinal_elem_implies_subseteq}.
                        Hence $q(\suc{\gamma})
                            \in\ran{\restrl{q}{\beta}}$ by
                            \auto{fun_ran_iff}.
                        Hence $q(\gamma)\in
                            \unions{\ran{\restrl{q}{\beta}}}$ by
                            \auto{unions_iff}.
                    \end{subproof}
                    Follows by \auto.
                \caseOf{$\beta$ is not a limit ordinal.}
                    Hence $\beta$ is a successor ordinal.
                    $\beta=\suc{\unions{\beta}}$ by
                        \obv{successor_ordinal_eq_suc_unions}.
                    Let $\eta=\unions{\beta}$.
                    $\eta\in\beta$ by
                        \auto{suc_intro_self}.
                    Hence $\eta\in\lambda$ by
                        \obv{ordinal_transitivity}.
                    $q(\suc{\eta})=
                        \normalDominatingNext{t}{\suc{\eta}}
                            {\restrl{q}{\suc{\eta}}}$ by
                        \auto{normal_dominating_recursion_successor}.
                    $\suc{\eta}\subseteq\lambda$ by
                        \obv{ordinal_elem_implies_subseteq}.
                    Hence $\suc{\eta}\subseteq\dom{q}$.
                    $\eta\in\suc{\eta}$ by \obv{suc_intro_self}.
                    $(\restrl{q}{\suc{\eta}})(\eta)=q(\eta)$ by
                        \obv{restrl_of_function_apply}.
                    $\unions{\suc{\eta}}=\eta$ by
                        \obv{ordinal_successor_data}.
                    Let $A=(\suc{\eta}\union q(\eta))
                        \union t(q(\eta))$.
                    $q(\suc{\eta})=\suc{A}$ by
                        \auto{normal_dominating_next}.
                    Hence $q(\beta)=\suc{A}$.
                    $q(\eta)\in\lambda$ by \obv{funs_type_apply}.
                    Hence $t(q(\eta))\in\lambda$ by
                        \obv{funs_type_apply}.
                    $\suc{\eta}$, $q(\eta)$, and $t(q(\eta))$
                        are ordinals by \obv{elem_of_ordinal_is_ordinal}.
                    Hence $A$ is an ordinal by
                        \obv{union_of_two_ordinals_is_ordinal}.
                    $\suc{A}$ is an ordinal by \obv{suc_ordinal}.
                    $\suc{\eta}\subseteq A$ by
                        \auto{union_upper_left,subseteq_transitive}.
                    Hence $\beta\subseteq q(\beta)$ by
                        \auto{subseteq_self_suc_intro,subseteq_transitive}.
                    Let $r=q(\eta)$.
                    $r$ is an ordinal.
                    $r\subseteq A$ by
                        \auto{union_upper_right,union_upper_left,
                            subseteq_transitive}.
                    Hence $r\in A$ or $r=A$ by
                        \obv{ordinal_subseteq_implies_elem_or_eq}.
                    Hence $r\in\suc{A}$ by \auto{suc_iff}.
                    Hence $q(\eta)\in\suc{A}$.
                    Hence $q(\eta)\in q(\beta)$.
                    For every $\xi\in\eta$ we have
                        $q(\xi)\in q(\eta)$ by \auto.
                    Show that for every $\gamma\in\beta$ we have
                        $q(\gamma)\in q(\beta)$.
                    \begin{subproof}
                        Fix $\gamma\in\beta$.
                        \begin{byCase}
                            \caseOf{$\gamma=\eta$.}
                                Follows.
                            \caseOf{$\gamma\neq\eta$.}
                                Hence $\gamma\in\eta$ by
                                    \auto{suc_iff}.
                                Hence $q(\gamma)\in q(\eta)$ by
                                    \auto.
                                Follows by \auto{ordinal_transitivity}.
                        \end{byCase}
                    \end{subproof}
                    Follows by \auto.
            \end{byCase}
    \end{byCase}
\end{proof}

\begin{theorem}\label{normal_dominating_recursion_is_normal_cofinal}
    Let $\lambda$ be a regular cardinal.
    Let $t\in\funs{\lambda}{\lambda}$.
    Let $q=\ordinalRecursion{\lambda}{\lambda}
        {\normalDominatingStep{t}{\lambda}}$.
    Then $\normalCofinalMap{q}{\lambda}{\lambda}$.
\end{theorem}
\begin{proof}
    Hence $\lambda$ is a limit ordinal by
        \obv{infinite_cardinal_is_limit_ordinal}.
    $q\in\funs{\lambda}{\lambda}$ by
        \auto{normal_dominating_recursion_type}.
    Hence $q$ is a function and $\dom{q}=\lambda$ by
        \obv{funs_is_function,funs}.
    $\ordinalMap{q}{\lambda}{\lambda}$ by \obv{ordinal_map}.
    Show that for all $j,k$ such that $j\in k\in\lambda$ we have
        $q(j)\in q(k)$.
    \begin{subproof}
        Fix $j,k$.
        Assume $j\in k\in\lambda$.
        Follows by \auto{normal_dominating_recursion_growth}.
    \end{subproof}
    Hence $\normalMapIncreasing{q}{\lambda}$ by
        \obv{normal_map_increasing}.
    Show that for every $l\in\lambda$ such that $l$ is a limit ordinal
        we have $q(l)=\unions{\ran{\restrl{q}{l}}}$.
    \begin{subproof}
        Fix $l\in\lambda$.
        Suppose $l$ is a limit ordinal.
        Follows by \auto{normal_dominating_recursion_limit}.
    \end{subproof}
    Hence $\normalMapContinuous{q}{\lambda}$ by
        \obv{normal_map_continuous}.
    Hence $\normalMap{q}{\lambda}{\lambda}$ by \obv{normal_map}.
    $\ran{q}\subseteq\lambda$ by \obv{funs_ran}.
    Show that for every $\alpha\in\lambda$ there exists
        $\gamma\in\ran{q}$ such that $\alpha\subseteq\gamma$.
    \begin{subproof}
        Fix $\alpha\in\lambda$.
        $q(\alpha)\in\ran{q}$ by \auto{fun_ran_iff}.
        $\alpha\subseteq q(\alpha)$ by
            \auto{normal_dominating_recursion_growth}.
        Follows by \auto.
    \end{subproof}
    Hence $\cofinalIn{\ran{q}}{\lambda}$ by \obv{cofinal_in}.
    Hence $\cofinalMap{q}{\lambda}{\lambda}$ by \obv{cofinal_map}.
    Follows by \obv{normal_cofinal_map}.
\end{proof}

\begin{theorem}\label{normal_dominating_recursion_later}
    Let $\lambda$ be a regular cardinal.
    Let $t\in\funs{\lambda}{\lambda}$.
    Let $\nu\in\lambda$ and $\mu\in\nu$.
    Let $q=\ordinalRecursion{\lambda}{\lambda}
        {\normalDominatingStep{t}{\lambda}}$.
    Then $t(q(\mu))\in q(\nu)$.
\end{theorem}
\begin{proof}
    Hence $\lambda$ is a limit ordinal by
        \obv{infinite_cardinal_is_limit_ordinal}.
    Hence $\lambda$ is an ordinal by \obv{limit_ordinal}.
    $\mu\in\lambda$ by \auto{ordinal_transitivity}.
    Hence $\suc{\mu}\in\lambda$ by \obv{suc_elem_limitordinal}.
    $q\in\funs{\lambda}{\lambda}$ by
        \auto{normal_dominating_recursion_type}.
    Hence $q(\nu)\in\lambda$ by \obv{funs_type_apply}.
    Hence $q(\nu)$ is an ordinal by
        \obv{elem_of_ordinal_is_ordinal}.
    $t(q(\mu))\in q(\suc{\mu})$ by
        \auto{normal_dominating_recursion_successor}.
    $\nu$ is an ordinal by \obv{elem_of_ordinal_is_ordinal}.
    Hence $\suc{\mu}\subseteq\nu$ by \obv{ordinal_suc_subseteq}.
    $\suc{\mu}$ is an ordinal by \obv{suc_ordinal}.
    Hence $\suc{\mu}\in\nu$ or $\suc{\mu}=\nu$ by
        \auto{ordinal_subseteq_implies_elem_or_eq}.
    \begin{byCase}
        \caseOf{$\suc{\mu}=\nu$.}
            Follows by \auto.
        \caseOf{$\suc{\mu}\neq\nu$.}
            Hence $\suc{\mu}\in\nu$ by \auto.
            Hence $q(\suc{\mu})\in q(\nu)$ by
                \auto{normal_dominating_recursion_growth}.
            Follows by \auto{ordinal_transitivity}.
    \end{byCase}
\end{proof}

\begin{theorem}\label{normal_dominating_map}
    Let $\lambda$ be a regular cardinal.
    Let $t\in\funs{\lambda}{\lambda}$.
    Then there exists $q$ such that
        $\normalCofinalMap{q}{\lambda}{\lambda}$ and
        for all $\mu,\nu$ such that $\mu\in\nu\in\lambda$ we have
            $t(q(\mu))\in q(\nu)$.
\end{theorem}
\begin{proof}
    Let $q=\ordinalRecursion{\lambda}{\lambda}
        {\normalDominatingStep{t}{\lambda}}$.
    $\normalCofinalMap{q}{\lambda}{\lambda}$ by
        \auto{normal_dominating_recursion_is_normal_cofinal}.
    For all $\mu,\nu$ such that $\mu\in\nu\in\lambda$ we have
        $t(q(\mu))\in q(\nu)$ by
        \auto{normal_dominating_recursion_later}.
    Follows by \auto.
\end{proof}
